\section{Modules (Done)}

\subsection{Exercise}\label{2.1} % 1

\begin{proof}
Let $x \otimes y \in (\Z/m\Z) \otimes_{\Z} (\Z/n\Z)$. Since $m, n$ are coprime, Bezout's lemma says that there are $a, b \in \Z$ such that $am + bn = 1$. Then,
$$x \otimes y = am(x \otimes y) + bn(x \otimes y) = (amx) \otimes y + x \otimes (bny) = 0 \otimes y + x \otimes 0 = 0,$$
so $(\Z/m\Z) \otimes_{\Z} (\Z/n\Z) = 0$.
\end{proof}

\subsection{Exercise}\label{2.2} % 2

\begin{proof}
Consider the exact sequence
$$0 \to \mathfrak{a} \to A \to A/\mathfrak{a} \to 0.$$
Tensoring with $M$, we have, by the right exactness of tensor products,
$$\mathfrak{a} \otimes_A M \to A \otimes_A M \to (A/\mathfrak{a})\otimes_A M \to 0.$$
Under canonical isomorphism, we have $A \otimes_A M = M$. Through this canonical isomorphism the image of $\mathfrak{a} \otimes_A M$ is generated via
$a \otimes m \mapsto am \in \mathfrak{a}M$. In fact, any element of $\mathfrak{a}M$ has the form $\sum_i a_im_i$ for $a_i \in \mathfrak{a}$ and
$m_i \in M$; this can be obtained through the map via the element $\sum_i(a_i,m_i) \in \mathfrak{a} \otimes_A M$. By exactness, we have
$$(A/\mathfrak{a})\otimes_A M \cong (A \otimes_A M)/\im(\mathfrak{a} \otimes_A M) \cong M / \mathfrak{a}M,$$
completing the proof.
\end{proof}

\subsection{Exercise}\label{2.3} % 3

\begin{proof}
Suppose $M \otimes_A N = 0$ for a local ring $A$ and finitely-generated $A$-modules $M$, $N$. If $k = A/\mathfrak{m}$ is the residue field,
let $M_k = k \otimes_A M = (A/\mathfrak{m}) \otimes_A M \cong M / \mathfrak{m}M$ (Exercise \ref{2.2}). Since $M\otimes_A N = 0$, we have
$$0 = k \otimes_A (M \otimes_A N) \otimes_A k = M_k \otimes_A N_k.$$
By Exercise \ref{2.2}, we have
$$M_k \otimes_A N_k \cong (M/\mathfrak{m}M) \otimes_A (N/\mathfrak{m}N).$$
Note that the $A$-linearity of $(M/\mathfrak{m}M) \otimes_A (N/\mathfrak{m}N)$ factors through $k = A/\mathfrak{m}$, so
$$0 = M_k \otimes_A N_k \cong M_k \otimes_k N_k.$$
However, note that since $M$ and $N$ are finitely-generated, $M_k$ and $N_k$ are finite-dimensional vector spaces over $k$. Thus, since the
tensor product vanishes, either $M_k$ or $N_k$ must vanish (dimensions multiply). By symmetry, assume $M_k = 0$ without loss of generality.

Since $A$ is a local ring, $\mathfrak{m}$ is identical to the Jacobson radical. Since $M/\mathfrak{m}M \cong M_k = 0$, we have that
$\mathfrak{m}M = M$. Since $M$ is also finitely-generated, we have by Nakayama's lemma that $M = 0$, as desired.
\end{proof}

\subsection{Exercise}\label{2.4} % 4

\begin{proof}
Suppose $M$ is flat and that we have an injective homomorphism of $A$-modules $j : N' \to N$. By flatness,
$$1_M \otimes j : M \otimes N' \to M \otimes N$$
is also injective. As $M \cong \bigoplus_i M_i$, we can define maps $1_{M_i} \otimes j : M_i \otimes N' \to M_i \otimes N$. Under canonical isomorphism, we have for
$N$ (and $N'$)
$$\bigoplus_{i \in I} (M_i \otimes N) \cong \left(\bigoplus_{i \in I} M_i\right) \otimes N \cong M \otimes N.$$
Thus, $1_M \otimes j = \oplus_{i\in I}(1_{M_{i}} \otimes j)$. By the flatness of $M$, we have
$$0 = \ker (1_M \otimes j) = \ker \left(\oplus_{i\in I} (1_{M_{i}} \otimes j)\right) = \oplus_{i\in I}\ker (1_{M_{i}} \otimes j),$$
so each $1_{M_{i}} \otimes j$ is injective, ensuring that $M_i$ is flat for all $i \in I$.

Conversely, suppose that $M_i$ is flat for all $i \in I$. Then, any exact sequence of $A$-modules
$$0 \to N' \to N \to N'' \to 0$$
yields the exact sequence
$$0 \to M_i \otimes N' \to M_i \otimes N \to M_i \otimes N'' \to 0$$
for all $i \in I$. Exact sequence commute with direct sums, so we have
$$0 \to \bigoplus_{i\in I}(M_i \otimes N') \to \bigoplus_{i\in I}(M_i \otimes N) \to \bigoplus_{i\in I}(M_i \otimes N'') \to 0$$
$$0 \to \left(\bigoplus_{i\in I}M_i\right) \otimes N' \to \left(\bigoplus_{i\in I}M_i\right) \otimes N \to \left(\bigoplus_{i\in I}M_i\right) \otimes N'' \to 0$$
$$0 \to M \otimes N' \to M \otimes N \to M \otimes N'' \to 0,$$
so $M$ is flat.
\end{proof}

\subsection{Exercise}\label{2.5} % 5

\begin{proof}
Consider $A[x]$ as an $A$-module. For any non-negative integers $n \ge 0$, we can view $M_n := Ax^n$ as an $A$-submodule since it is additively-closed, and $A$ has a closed action on it. 
As a module, $M_n \cong A$ under the $A$-module homomorphism $a \mapsto ax^n$. Thus, for any $A$-module $N$, $M_n \otimes_A N \cong N$ under $M_n \cong A$ and the canonical isomorphism 
$A \otimes_A N \cong N \cong N$. Thus, tensoring a short exact sequence of $A$-modules with $M_n$ yields an isomorphic short exact sequence, so by definition, $M_n$ is a flat $A$-module.

Now, notice that any $f \in A[x]$ can be written as a finite and unique sum of elements in various $M_n$. In particular, $A[x] \cong \bigoplus_{n \ge 0} M_n$ as $A$-modules. 
Invoking Exercise \ref{2.4}, $A[x]$ is also flat.
\end{proof}

\subsection{Exercise}\label{2.6} % 6

\begin{proof}
If $M$ is an $A$-module, then we can define an action of $A[x]$ on $M[x]$ in a natural way. More explicitly, the action of $ax^i \in A[x]$ on $mx^j \in M[x]$ is
$(am)x^{i+j} \in M[x]$, exploiting the fact that $am \in M$. We extend the action $A$-linearly to all of $A[x]$ and $M[x]$, so $M[x]$ has an $A[x]$-action. Clearly, $M[x]$ is an 
abelian group additively. 
Thus, the only non-trivial condition to check on $M[x]$ being an $A[x]$-module is $(fg)m = f(gm)$. Thus, let $f = \sum_i f_i x^i$, $g = \sum_j g_j x^j$, and $m = \sum_k m_k x^k$, 
where $m \in M[x]$ and $f,g \in A[x]$. Then,
$$(fg)m = \left(\sum_i\sum_j f_ig_j x^{i+j}\right)\left(\sum_k m_k x^k\right) = \sum_i\sum_j\sum_k  f_ig_jm_k x^{i+j+k}$$
and
$$f(gm) = \left(\sum_i f_ix^i\right)\left(\sum_j\sum_k g_jm_k x^{j+k}\right) = \sum_i\sum_j\sum_k  f_ig_jm_k x^{i+j+k}$$
so $(fg)m = f(gm)$. Thus, $M[x]$ is an $A[x]$-module.

Now, we show that $M[x] \cong A[x] \otimes_A M$ as $A[x]$-modules. We use the direct product decomposition of Exercise \ref{2.5}, where as an $A$-module, $A[x] = \bigoplus_{i\ge 0}A x^i$; we saw that
$A \cong Ax^i$ as an $A$-module via $a \mapsto ax^i$. Thus, under canonical isomorphism, we have
$$A[x] \otimes_A M \cong \left(\bigoplus_{i\ge 0}A x^i\right)\otimes_A M \cong \bigoplus_{i\ge 0}\left((A x^i)\otimes_A M\right).$$
Since $Ax^i \cong A$ and $A \otimes_A M \cong M$, both as $A$-modules, we can compose the maps to obtain the isomorphism of $A$-modules $(Ax^i)\otimes_A M \cong Mx^i$. Thus, we have
$$A[x] \otimes_A M \cong \bigoplus_{i\ge 0}\left((A x^i)\otimes_A M\right) \cong \bigoplus_{i\ge 0} Mx^i \cong M[x].$$
As $A$-modules then, $A[x]\otimes_A M\cong M[x]$. This establishes a bijection. To show this is truly an $A[x]$-module homomorphism, we'll show the action commutes with $(x^i,m) \mapsto mx^i$; terms
of this form generate the modules $A$-linearly, so this will suffice. Let $f = \sum_j f_j x^j \in A[x]$. Then,
$$f(x).(x^i\otimes m) = (f(x)\otimes 1)(x^i\otimes m) = \left(\sum_j f_j x^{i+j}\right)\otimes m \\$$
$$\mapsto \sum_j f_jm x^{i+j} = \left(\sum_j f_j x^j\right)(mx^i) = f(x).(mx^i),$$
so we have the desired homomorphism. Thus, $M[x] \cong A[x] \otimes_A M$ as $A[x]$-modules.
\end{proof}

\subsection{Exercise}\label{2.7} % 7

\begin{proof}
Let $\mathfrak{p}$ be a prime ideal of $A$. Since $\mathfrak{p}$ is an $A$-module, by Exercise \ref{2.6}, $\mathfrak{p}[x]$ is an $A[x]$-module, so it is an ideal of $A[x]$. Suppose $f,g \in A[x]$ 
and $fg \in \mathfrak{p}[x]$. With the map $A \to A/\mathfrak{p}$, we have a surjective $A[x]$-module homomorphism $A[x] \to (A/\mathfrak{p})[x]$. Under this map, $\overline{fg} = 0$. Suppose that 
$\overline{f}$ is a zero-divisor. Then, by Exercise \ref{1.2}, there is a nonzero element $a \in A/\mathfrak{p}$ such that $a\overline{f} = 0$. Thus, $\sum_i af_i x^i = 0$, but since $A/\mathfrak{p}$ 
is an integral domain, we must have $f_i = 0$ for all $i$, i.e. $\overline{f} = 0$. Hence, $f \in \mathfrak{p}[x]$, and $\mathfrak{p}[x]$ is a prime ideal of $A[x]$.
\end{proof}

\begin{remark}
Note that if $\mathfrak{m}$ is a maximal ideal of $A$, then $\mathfrak{m}[x]$ is not maximal as an ideal of $A[x]$. Simply note that $\mathfrak{m} \subset \mathfrak{m} + xA[x] \subset A[x]$
is a proper chain of ideals.
\end{remark}

\subsection{Exercise}\label{2.8} % 8

\begin{proof}[(i) Proof.]
Suppose $M$ and $N$ are flat $A$-modules. Let $j : J \to J'$ be any injection of $A$-modules. Then, since $N$ is flat,
$$(1 \otimes j) : N \otimes J \to N \otimes J'$$
is also injective. But $M$ is also flat, so
$$1 \otimes (1 \otimes j) : M \otimes (N \otimes J) \to M \otimes (N \otimes J')$$
is injective. Recall that we have the natural isomorphism $M \otimes (N \otimes J) \cong (M \otimes N) \otimes J$ (same with $J'$ in place of $J$). Thus, we have an injection
$$1 \otimes j : (M \otimes N) \otimes J \to (M \otimes N) \otimes J'.$$
Hence, $M \otimes_A N$ is a flat $A$-module.
\end{proof}

\begin{proof}[(ii) Proof.]
Suppose $B$ is a flat $A$-algebra and $N$ is a flat $B$-module. Let $j : J \to J'$ be any injection of $A$-modules. Then, since $B$ is flat as an $A$-module,
$$(1 \otimes j) : B \otimes_A J \to B \otimes_A J'$$
is injective. Also, note that $B \otimes_A J$ and $B \otimes_A J'$ are naturally $B$-modules. Since $N$ is flat as a $B$-module,
$$1 \otimes (1 \otimes j) : N \otimes_B (B \otimes_A J) \to N \otimes_B (B \otimes_A J')$$
is injective. Since $B$ is an $(A,B)$-bimodule, under canonical isomorphism, we have an injection
$$1 \otimes j : (N \otimes_B B) \otimes_A J \to (N \otimes_B B) \otimes_A J'.$$
Yet again, under canonical isomorphism, we have an injection
$$1 \otimes j : N \otimes_A J \to N \otimes_A J',$$
so $N$ is flat as an $A$-module.
\end{proof}

\subsection{Exercise}\label{2.9} % 9

\begin{proof}
Let
$$0 \to M' \xlongrightarrow{f} M \xlongrightarrow{g} M'' \to 0$$
be an exact sequence of $A$-modules, where $M'$ and $M''$ are finitely-generated. By exactness, we have the isomorphism
$$M/f(M') \cong M''.$$
Since $M''$ is finitely-generated, $M/f(M')$ is generated by finitely-many elements $\{y_i\} \subset M$. Likewise, since $f$ is injective and $M'$ is finitely-generated, $f(M')$ is 
generated by finitely-many elements $\{x_j\} \subset M$. We claim $\{y_i\} \cup \{x_j\}$ generate $M$, so suppose $m \in M$. Then, 
$$m + f(M') = \sum_i a_i y_i + f(M')$$ 
for $a_i \in A$, so $m - \sum_i a_i y_i \in f(M')$. Thus,
$$m - \sum_i a_i y_i = \sum_j b_j x_j$$
for $b_j \in A$, so $m = \sum_i a_i y_i + \sum_j b_j x_j$. Hence, $M$ is finitely-generated.
\end{proof}

\subsection{Exercise}\label{2.10} % 10

\begin{proof}
By assumption, $u$ induces a surjective homomorphism $M/\mathfrak{a}M \to N/\mathfrak{a}N$, so $N = \mathfrak{a}N + \im u$. Note that $\im u$ is an $A$-submodule of $N$.
Thus, since $N$ is a finitely-generated $A$-module and $\mathfrak{a}$ is in the Jacobson radical of $A$, we have by Corollary 2.7 that
$N = \im u$, so $u$ is surjective.
\end{proof}

\subsection{Exercise}\label{2.11} % 11

\begin{proof}
Let $A \neq 0$ be a ring, so it has a maximal ideal $\mathfrak{m}$. Suppose $\phi : A^m \to A^n$ is an $A$-module isomorphism. Then, if $k = A/\mathfrak{m}$ is the residue field, we have an
isomorphism
$$1 \otimes \phi : k \otimes A^m \to k \otimes A^n.$$
Note that for $N > 0$, $A^N = \oplus_{i=1}^N A$, so $k \otimes A^N \cong \oplus_{i=1}^N(k \otimes A)$ by natural isomorphism, and $k \otimes_A A \cong k$ by natural isomorphism. Thus,
our map $1 \otimes \phi$ is a map $k^m \to k^n$. This is an $A$-module isomorphism and thus a $k$-module (vector space) isomorphism. Hence, $m = n$.
\end{proof}

\begin{proof}
Now, suppose $A \neq 0$, and $\phi : A^m \to A^n$ is a surjective $A$-module homomorphism. Since $A$ is nonzero, it has a maximal ideal $\mathfrak{m}$, so let $k = A/\mathfrak{m}$ be the
residue field. Then, we have a surjection (right exactness)
$$1 \otimes \phi : k \otimes A^m \to k \otimes A^n.$$
Perfectly analogous to the last proof, $k \otimes A^N \cong k^N$, so we have a surjective homomorphism of $A$-modules $k^m \to k^n$. Thus, it is also a surjective $k$-module (vector space) homomorphism,
so $m \ge n$.
\end{proof}

\begin{proof}
Again, assume $A \neq 0$ is a ring. Here, we prove that if $\phi : A^m \to A^n$ is injective, then $n \ge m$. We'll prove this by contradiction. Suppose we have an injective $A$-module homomorphism $\phi: A^m \to A^n$ with $m > n$. We can
compose $\phi$ with the map $\psi: A^n \to A^n \oplus 0^{m-n} \subset A^m$ so that $\epsilon = \psi \circ \phi$ is an endomorphism of the $A$-module $A^m$. As the composition of two injective maps, it is itself injective.
Note that $\epsilon(A^m) \subset A^n \oplus 0^{m-n}$. 
Since $A^m$ is clearly finitely-generated as an $A$-module, by Proposition 2.4, $\epsilon$ satisfies an equation of the form:
$$\epsilon^k + a_1 \epsilon^{k-1} + \dots + a_k \id_{A^m} = 0,$$
where $a_i \in A$ for all $i$. Choose an equation of this form such that $k$ is minimal. Note that if $k = 0$, then $0 = \id_{A^m}$, but $A \neq 0$, so $k > 0$.

Consider the projection map $\pi : A^m \to A$ that projects onto the last direct summand. Since $m > n$ by assumption, $\pi \circ \epsilon = 0$. Thus, taking $\pi$ on both sides, 
we have $\pi(a_k \id_{A^m}) = 0$. But, 
$$\pi(a_k \id_{A^m}((0,\ldots,0,1))) = \pi(0,\ldots,0,a_k) = a_k,$$ 
so $a_k = 0$. Thus, our equation is
$$\epsilon \circ (\epsilon^{k-1} + a_1 \epsilon^{k-2} + \dots + a_{k-1} \id_{A^m}) = 0.$$
Since $\epsilon$ is injective,
$$\epsilon^{k-1} + a_1 \epsilon^{k-2} + \dots + a_{k-1} \id_{A^m} = 0,$$
violating the minimality of $k$ and giving us the desired contradiction.
\end{proof}

\subsection{Exercise}\label{2.12} % 12

\begin{proof}
Let $\{e_1,\ldots,e_n\}$ be an $A$-basis of $A^n$, and choose $u_i \in M$ such that $\phi(u_i) = e_i$. This can be done since the given $A$-module homomorphism $M \to A^n$ is surjective. 
Let $m \in M$, and let $\sum_i a_i e_i$ be the unique decomposition of $\phi(m) \in A^n$, where $a_i \in A$. Then, $\phi\left(m - \sum_i a_i u_i\right) = 0$, so $m - \sum_i a_i u_i \in \ker\phi$.
Also, if $\sum_i a_i u_i \in \ker \phi$, then $0 = \phi\left(\sum_i a_i u_i\right) = \sum_i a_i e_i$, but since $\{e_i\}$ is an $A$-basis, $a_i = 0$ for all $i$. If $U$ is the $A$-submodule generated
by $\{u_i\}$, we've shown that $\ker \phi \cap U = \{0\}$ and $M = \ker \phi + U$. Hence, $M = \ker \phi \oplus U$.

Since $M$ is finitely-generated, we have finitely-many generators $\{v_j\}$ of $M$. Under our direct sum decomposition, we have a projection map $\pi : M \to \ker \phi$. 
Thus, $\{\pi(v_j)\}$ generate $\ker \phi$, so $\ker\phi$ is finitely-generated.
\end{proof}

\subsection{Exercise}\label{2.13} % 13

\begin{proof}
If $f : A \to B$ is a ring homomorphism and $N$ is a $B$-module, we can restrict scalars to $A$, then extend scalars back to $B$. Thus, we have another $B$-module
$N_B := B \otimes_A N$ ($b \in B$ acts via $b \otimes 1$) and a $B$-module homomorphism $g: N \to N_B$ via $y \mapsto 1 \otimes y$. We have a canonical ``absorption" map $p : N_B \to N$ that maps
$b \otimes y \mapsto by$. Let $1 \otimes y \in \im g$; $\im g$ is generated by elements of this form. Then, if $1 \otimes y \in \ker p$, $0 = p(1\otimes y) = y$, so $1 \otimes y = 0$. 
Thus, $\im g \cap \ker p = \{0\}$. Consider $b \otimes y \in N_B$; elements of this form generate $N_B$. Note that $b \otimes y = (b \otimes y - 1 \otimes (by)) + 1 \otimes (by)$,
and that $b \otimes y - 1 \otimes (by) \in \ker p$ and $1 \otimes (by) \in \im g$. Thus, $N_B = \im g \oplus \ker p$.
\end{proof}

\subsection{Exercise}\label{2.14} % 14

A partially ordered set $I$ is said to be a {\it directed} set if for each pair $i,j \in I$ there exists $k \in I$ such that $i \le k$ and $j \le k$.

Let $A$ be a ring, let $I$ be a directed set, and let $(M_i)_{i \in I}$ be a family of $A$-modules indexed by $I$. For each pair $i,j \in I$ such that
$i \le j$, let $\mu_{ij} : M_i \to M_j$ be an $A$-homomorphism, and suppose that the following axioms are satisfied:
\begin{itemize}
    \item[(1)] $\mu_{ii}$ is the identity mapping of $M_i$ for all $i \in I$;
    \item[(2)] $\mu_{ik} = \mu_{jk} \circ \mu_{ij}$ whenever $i \le j \le k$.
\end{itemize}
Then, the modules $M_i$ and homomorphisms $\mu_{ij}$ are said to form a {\it direct system} $\mathbf{M} = (M_i, \mu_{ij})$ over the directed set $I$.

We shall construct an $A$-module $M$ called the {\it direct limit} of the direct system $\mathbf{M}$. Let $C$ be the direct sum of the $M_i$ and identity
each module $M_i$ with its canonical image in $C$. Let $D$ be the submodule of $C$ generated by all elements of the form $x_i - \mu_{ij}(x_i)$ where $i \le j$
and $x_i \in M_i$. Let $M = C/D$, let $\mu : C \to M$ be the projection, and let $\mu_i$ be the restriction of $\mu$ to $M_i$.

The module $M$, or more correctly the pair consisting of $M$ and the family of homomorphisms $\mu_i : M_i \to M$, is called the {\it direct limit} of the direct
system $\mathbf{M}$, and is written $\varinjlim M_i$. For $x_i \in M_i$ and $i\le j$, we have by construction that $x_i - \mu_{ij}(x_i) \in D = \ker \mu$.
Thus, $\mu_i = \mu_j \circ \mu_{ij}$ since
$$\mu_i(x_i) = \mu(x_i) = \mu(\mu_{ij}(x_i)) = \mu_j(\mu_{ij}(x_i)).$$

\subsection{Exercise}\label{2.15} % 15

\begin{proof}
In the construction of $M$ in Exercise \ref{2.14}, we saw that $M$ is a quotient of the direct sum $C = \bigoplus_{i \in I}M_i$. Let $y \in M$, so there is some $x \in C$ such that $\mu(x) = y$. 
Then, $x = \sum_{j \in J}x_j$ for a finite subset $J \subset I$. Since $I$ is a directed set, for any two elements $i,j \in I$, there exists a $k \in I$ such that $i \le k$ and $j \le k$. 
Since $J$ is a finite subset, a routine inductive argument shows there is some $i \in I$ such that $j \le i$ for all $j \in J$. Thus, 
$$y = \mu(x) = \mu\left(\sum_j x_j\right) = \sum_j \mu(x_j) = \sum_j \mu_j(x_j) = \sum_j (\mu_i\circ\mu_{ji})(x_j) = \mu_i\left(\sum_j \mu_{ji}(x_j)\right),$$
where we invoked $\mu_j = \mu_i \circ \mu_{ji}$ whenever $j \le i$ from Exercise \ref{2.14}. Thus, letting $x_i = \sum_j \mu_{ji}(x_j)$, we see that $y = \mu_i(x_i)$ for some $i \in I$ and 
$x_i \in M_i$, as desired.
\end{proof}

\begin{proof}
Let's first lay some groundwork. $D$, the kernel of $\mu$, is generated by elements of the form $x_i - \mu_{ij}(x_i)$. In other words, the general element has the form
$$\sum_k a_k(x_{i_k} - \mu_{i_k j_k}(x_{i_k})).$$
But, since $\mu_{**}$ is an $A$-module homomorphism, we have $a_k(x_{i_k} - \mu_{i_k j_k}(x_{i_k})) = (a_k x_{i_k}) - \mu_{i_k j_k}(a_k x_{i_k})$, so we can always choose a representation
such that $a_k = 1$ for all $k$. Further, if we have a repeated pair $(i,j) = (i_k,j_k)$, then $(x_1 - \mu_{ij}(x_1)) + (x_2 + \mu_{ij}(x_2)) = (x_1 + x_2) - \mu_{ij}(x_1 + x_2)$, so we can
choose a representation such that pairs $(i,j)$ occur only once.

Suppose $\mu_i(x_i) = 0$. Then, $x_i \in M_i \cap D$, so it has a representation
$$x_i = \sum_{(j,k) \in P}(y_j - \mu_{jk}(y_j)) = \sum_{j \in J} z_j$$
for a finite subset $J \subset I$. Since we are currently in the direct sum $C$ and $x_i \in M_i$, the projection map onto $M_j$ for $i \neq j$ applied to the equation vanishes. Since the right-hand side breaks the middle
sum into $M_j$ components, we have $z_j = 0$ if $j \neq i$ and $z_i = x_i$. Let $\ell \in I$ be such that $j \le \ell$ for all $j\in J$; such an element can be found since $I$ is a directed set and $J$
is finite. Since $z_i = x_i$ and $z_j = 0$ for $j \neq i$, we have trivially that $\mu_{i\ell}(x_i) = \sum_{j \in J}\mu_{j\ell}(z_j)$. Thus,
$$\mu_{i\ell}(x_i) = \sum_{j \in J}\mu_{j\ell}(z_j) = \sum_{(j,k) \in P}(\mu_{j\ell}(y_j) - (\mu_{k\ell}\circ\mu_{jk})(y_j)) 
= 0$$
since $\mu_{j\ell} = \mu_{k\ell}\circ\mu_{jk}$, as desired.
\end{proof}

\subsection{Exercise}\label{2.16} % 16

\begin{proof}
Let's start by showing the direct limit $M$ satisfies these properties. Let $N$ be an $A$-module such that for each $i \in I$, $\alpha_i : M_i \to N$ is an $A$-module homomorphism
such that $\alpha_i = \alpha_j \circ \mu_{ij}$ whenever $i \le j$. Let $y \in M$. By Exercise \ref{2.15}, $y = \mu_i(x_i)$ for some $x_i \in M_i$. We'd like to define $\alpha : M \to N$
such that $\alpha(y) = \alpha_i(x_i)$, but we have to first show this is well-defined. Let $x_i' \in M_i$ be another element such that $y = \mu_i(x_i')$. Then, $\mu_i(x_i - x_i') = 0$. So,
to show the map $\alpha$ is well-defined, we have to show that $\mu_i(y_i) = 0$ implies $\alpha_i(y_i) = 0$ (and thus that $\alpha_i(x_i) = \alpha_i(x_i')$). By Exercise \ref{2.15}, we know that $\mu_i(y_i) = 0$ 
implies there is some $j \ge i$ such that $\mu_{ij}(y_i) = 0$. By assumption, $\alpha_i(y_i) = \alpha_j(\mu_{ij}(y_i)) = \alpha_j(0) = 0$, so $\alpha$ is well-defined. By definition, $\alpha_i = \alpha \circ \mu_i$.
Also, this homomorphism is clearly unique since if $y \in M$, $y = \mu_i(x_i)$ for some $i \in I$ and $x_i \in M_i$, so $\alpha(y) = \alpha(\mu_i(x_i)) = \alpha_i(x_i)$. In particular, $\alpha$
is completely determined by the given data $\{(M_i,\alpha_i)\}_{i \in I}$.

We now show uniqueness. Let $M'$ satisfy the given properties. Then, if $M$ is the direct limit, $M$ is an $A$-module with compatible maps $\mu_i : M_i \to M$ since $\mu_i = \mu_j \circ \mu_{ij}$ for 
$j \ge i$ by Exercise \ref{2.14}. Then, $\mu_i = \alpha' \circ \mu_i'$ for all $i$, where $\alpha' : M' \to M$ is an $A$-module homomorphism. Symmetrically, we have a map $\alpha : M \to M'$
such that $\mu_i' = \alpha \circ \mu_i$. Thus, we have $(\alpha' \circ \alpha)$ is the identity on the image of $M_i$ in $M$ for all $i$; every element of $M$ has this form, so $\alpha' \circ \alpha = \id_M$.
Similarly, $\alpha \circ \alpha' = \id_{M'}$, so $\alpha : M \to M'$ is an $A$-module isomorphism.
\end{proof}

\subsection{Exercise}\label{2.17} % 17

\begin{proof}
By construction, $\varinjlim M_i = \left(\bigoplus_{i\in I} M_i\right)/D$, where $D$ is the $A$-module generated by $x_i - \mu_{ij}(x_i)$ for $i \le j$. Consider the canonical homomorphism
$\psi : \left(\bigoplus_{i\in I} M_i\right) \to \sum M_i$. In this case, $\mu_{ij}$ is the map that identifies $M_i$ as a submodule of $M_j$. In particular, the submodule $D$ of the direct sum 
is in $\ker\psi$ since $x_i - \mu_{ij}(x_i) = 0 \in \sum_k M_k$ for $i \le j$. Thus, we have a well-defined $A$-module homomorphism
$$\phi : \left(\bigoplus_{i\in I} M_i\right)/D \to \sum M_i.$$
It is clearly surjective, since $\sum_i x_i \in \sum M_i$ satisfies $\psi\left(\sum_i x_i\right) = \sum_i x_i$. Now, suppose we have $\phi\left(D + \sum_i x_i\right) = 0$, i.e. $\sum_i x_i = 0$ in $\sum M_i$.
Since there are finitely-many $i$ indexing this sum, there is some $k$ such that for each $i$, $x_i \in M_k$. Then,
$$D + \sum_i x_i = \left(D + \sum_i(\mu_{ik}(x_i) - x_i)\right) + \sum_i x_i = D + \sum_i\mu_{ik}(x_i).$$
In $\bigoplus_{i \in I}M_i$, $\sum_i\mu_{ik}(x_i)$ lies entirely in the direct summand $M_k$; since it is $0$ in $M_k$, it is zero in the direct sum, so $\phi$ is injective. Thus, $\phi$ is an
$A$-module isomorphism, so $\varinjlim M_i \cong \sum M_i$. Any element of $\sum M_i$ is a finite sum $\sum_i x_i$, where $x_i \in M_i$. Since $I$ is directed, there is some $M_k$ that contains each $M_i$, 
so $\sum M_i \subset \bigcup M_i$. The converse clearly holds, so $\varinjlim M_i \cong \sum M_i \cong \bigcup M_i$.
\end{proof}

\subsection{Exercise}\label{2.18} % 18

\begin{proof}
Consider the $A$-module homomorphisms $\alpha_i = \nu_i \circ \phi_i$ from $M_i$ to $N$. If $i \le j$, then
$$\alpha_j \circ \mu_{ij} = \nu_j \circ \phi_j \circ \mu_{ij} = \nu_j \circ \nu_{ij} \circ \phi_i = \nu_i \circ \phi_i = \alpha_i,$$
so by the universal mapping property of direct limits (Exercise \ref{2.16}), we have a unique homomorphism $\phi : M \to N$ such that $\alpha_i = \phi \circ \mu_i$
for all $i$. Unwinding the definition, $\phi \circ \mu_i = \nu_i \circ \phi_i$ for all $i \in I$, as desired.
\end{proof}

\subsection{Exercise}\label{2.19} % 19

\begin{proof}
Suppose we have directed systems of $A$-modules $\mathbf{M}$, $\mathbf{N}$, and $\mathbf{P}$ over the same directed set $I$ such that
$$M_i \xlongrightarrow{f_i} N_i \xlongrightarrow{g_i} P_i$$
is exact for all $i \in I$. Consider the sequence of direct limits
$$M \xlongrightarrow{f} N \xlongrightarrow{g} P,$$
where $f = \varinjlim f_i$ and $g = \varinjlim g_i$ (Exercise \ref{2.18}). Let $\mu_{ij}$, $\nu_{ij}$, and $\pi_{ij}$ be the corresponding
compatibility maps. Suppose $m \in M$. Then, by Exercise \ref{2.15}, there is some $i \in I$ such that $m = \mu_i(x_i)$. Then,
$$(g \circ f)(m) = (g \circ f \circ \mu_i)(x_i) = (g \circ \nu_i \circ f_i)(x_i) = \pi_i((g_i \circ f_i)(x_i)) = \pi_i(0) = 0$$
since $g_i \circ f_i = 0$ by exactness. Thus, $g \circ f = 0$, so $\im f \subset \ker g$.

Now, let $n \in \ker g$. By Exercise \ref{2.15}, there is some $i \in I$ such that $n = \nu_i(x_i)$. Thus,
$$0 = g(n) = (g \circ \nu_i)(x_i) = (\pi_i \circ g_i)(x_i) = \pi_i(g_i(x_i)).$$
By Exercise \ref{2.15}, there is some $j \ge i$ such that $\pi_{ij}(g_i(x_i)) = 0$. So,
$$0 = (\pi_{ij} \circ g_i)(x_i) = (g_j \circ \nu_{ij})(x_i) = g_j(\nu_{ij}(x_i)) = g_j(x_j),$$
where we define $x_j = \nu_{ij}(x_i)$. Note that $n = \nu_i(x_i) = \nu_j(\nu_{ij}(x_i)) = \nu_j(x_j)$ and $x_j \in \ker g_j$. By exactness, $x_j \in \im f_j$, so there
is some $y_j \in M_j$ such that $x_j = f_j(y_j)$. Let $y = \mu_j(y_j)$. Then,
$$f(y) = (f \circ \mu_j)(y_j) = (\nu_j \circ f_j)(y_j) = \nu_j(f_j(y_j)) = \nu_j(x_j) = n,$$
so $n \in \im f$. Thus, $\ker g \subset \im f$, so finally, $\ker g = \im f$, giving us the desired exactness.
\end{proof}

\subsection{Exercise}\label{2.20} % 20

\begin{proof}
Let $N$ be any $A$-module. Then, $(M_i \otimes N, \mu_{ij} \otimes 1)$ is a direct system; let $P = \varinjlim(M_i \otimes N)$. For
each $i \in I$ we have a homomorphism $\mu_i \otimes 1 : M_i \otimes N \to M \otimes N$, so by Exercise \ref{2.16}, we have a homomorphism
$\psi : P \to M \otimes N$.

For each $i \in I$, define $g_i : M_i \times N \to M_i \otimes N$ to be the canonical $A$-bilinear map. For every $n \in N$, this defines an
$A$-module homomorphism $g_i(-,n) \to M_i \otimes N$. Passing to the direct limit, we obtain $g_n = \varinjlim g_i(-,n)$, which satisfies
$g_n \circ \mu_i = (\mu_i \otimes 1)(g_i(-,n))$. We define $g : M \times N \to P$ as $g(m,n) = g_n(m)$. This is $A$-linear in the first coordinate, so we must show 
it in the second. 
$$(g_n + g_m)\circ \mu_i = g_n \circ \mu_i + g_m \circ \mu_i = (\mu_i \otimes 1)(g_i(-,n)) + (\mu_i \otimes 1)(g_i(-,m))$$
$$= (\mu_i \otimes 1)(g_i(-,n+m)) = g_{n+m}\circ \mu_i$$
for all $i$. By the uniqueness of the direct limit map (Exercise \ref{2.18}), $g_n + g_m = g_{n+m}$. Also,
$$a(g_n \circ \mu_i) = a (\mu_i \otimes 1)(g_i(-,n)) = (\mu_i \otimes 1)(g_i(-,an)) = g_{an} \circ \mu_i$$
for all $i$. By the uniqueness of the direct limit map again, $a g_n = g_{an}$, so $g : M \times N \to P$ is $A$-bilinear. Thus, $g$ factors through $M \otimes N$
giving us an $A$-module homomorphism
$$\phi : M \otimes N \to P.$$

Consider the map $\psi \circ \phi$. Since it is an $A$-module endomorphism on $M \otimes N$, it is defined by its image on generators $m \otimes n$. Since $m \in M$, note
that there is some $i \in I$ such that $\mu_i(x_i) = m$ for $x_i \in M_i$. Then,
$$(\psi \circ \phi)(m \otimes n) = \psi(g(m,n)) = \psi(g_n(m)) = \psi((g_n \circ \mu_i)(x_i)) = \psi(g_i(x_i,n))$$ 
$$= \psi(x_i \otimes n) = (\mu_i \otimes 1)(x_i \otimes n) = \mu_i(x_i) \otimes n = m \otimes n,$$
so $\psi \circ \phi = 1_{M\otimes N}$. Conversely, let $p \in P$. Then, by Exercise \ref{2.15}, there is some $i \in I$ and $\sum x_i \otimes n \in M_i \otimes N$ such that
$p = (\mu_i \otimes 1)\left(\sum x_i \otimes n \right)$. Since we already know that $\phi \circ \psi$ is an $A$-module homomorphism, it is enough to show that
$(\phi \circ \psi)((\mu_i\otimes 1)(x_i \otimes n)) = (\mu_i \otimes 1)(x_i \otimes n)$ to see that $\phi \circ \psi = \id_P$. So let's prove that:
$$(\phi \circ \psi)((\mu_i\otimes 1)(x_i \otimes n)) = \phi(\mu_i(x_i) \otimes n) = g(\mu_i(x_i),n)$$
$$= g_n(\mu_i(x_i)) = (\mu_i \otimes 1)(g_i(x_i,n)) = (\mu_i \otimes 1)(x_i \otimes n),$$
as desired.
\end{proof}

\subsection{Exercise}\label{2.21} % 21

\begin{proof}
Let $\{(A_i,\alpha_{ij})\}$ be a directed system of rings so that $\alpha_{ij} : A_i \to A_j$ is a ring homomorphism; trivially, they are $\Z$-modules and $\alpha_{ij}$
are $\Z$-module homomorphisms, so via the construction of Exercise \ref{2.14}, we have a direct limit $A = \varinjlim A_i$. $A$ is thus a $\Z$-module, so $(A,+)$ is an
abelian group.

Let $x,y \in A$. Then, by Exercise \ref{2.15}, there are $i,j \in I$, $x_i \in A_i$, and $y_j \in A_j$ such that $\alpha_i(x_i) = x$ and $\alpha_j(y_j) = y$. Since
$I$ is a directed set, there is some $k$ such that $i \le k$ and $j \le k$. Consider then, $x_k = \alpha_{ik}(x_i)$ and $y_k = \alpha_{jk}(y_j)$. Note that
$$\alpha_k(x_k) = (\alpha_k \circ \alpha_{ik})(x_i) = \alpha_i(x_i) = x,$$
and similiarly, $\alpha_k(y_k) = y$. Thus, we define $x \cdot y = \alpha_k(x_k \cdot y_k)$, using the ring structure of $A_k$. To show that this is well-defined, we have
to show that this is independent of both $k$ and the representatives $x_k$ and $y_k$. Suppose we have another representative $x_k'$ of $x$. Then, 
$$\alpha_k(x_k' \cdot y_k) - \alpha_k(x_k \cdot y_k) = \alpha_k((x_k' - x_k) \cdot y_k).$$
Note that since $\alpha_k(x_k' - x_k) = 0$, we have by Exercise \ref{2.15}, that there is some $\ell \ge k$ such that $\alpha_{k\ell}(x_k' - x_k) = 0$. Since
$\alpha_{k\ell}$ is a ring homomorphism, we have
$$\alpha_k((x_k' - x_k) \cdot y_k) = (\alpha_{\ell}\circ\alpha_{k\ell})((x_k' - x_k) \cdot y_k) = \alpha_{\ell}(\alpha_{k\ell}(x_k' - x_k)\cdot\alpha_{k\ell}(y_k)) = \alpha_{\ell}(0\cdot\alpha_{k\ell}(y_k)) = 0,$$
By exactly the same argument, or rather by symmetry, if $y_k'$ is another representative of $y$, the same invariance holds: i.e.,
$$\alpha(x_k \cdot y_k) = \alpha(x_k' \cdot y_k) = \alpha(x_k' \cdot y_k').$$
Thus, we've shown the multiplication does not depend on the particular representations of $x$ and $y$. 

Now, we must show that it does not depend on $k$. Let $i \in I$ be such that there are $x_i, y_i \in A_i$ satisfying 
$x = \alpha_i(x_i)$ and $y = \alpha_i(y_i)$. Then, since $I$ is directed, there is some $j \in J$ such that $i \le j$ and $k \le j$. Define $x_{ij} = \alpha_{ij}(x_i)$, and define $x_{kj}$, $y_{ij}$, and $y_{kj}$ similarly.
Then, 
$$\alpha_j(x_{ij}) = \alpha_{j}(\alpha_{ij}(x_i)) = \alpha_i(x_i) = x;$$
the same argument applies to $x_{kj}$, $y_{ij}$, and $y_{kj}$. In particular, $x_{ij}$ and $x_{kj}$ are two representations of $x$ in $A_j$, and $y_{ij}$ and $y_{kj}$ are two representations of $y$
in $A_j$. Since we've already shown the representations in the same ring produce the same multiplication, we have 
$$\alpha_i(x_i \cdot y_i) = \alpha_j(\alpha_{ij}(x_i \cdot y_i)) = \alpha_j(x_{ij} \cdot y_{ij}) = \alpha_j(x_{kj} \cdot y_{kj}) = \alpha_j(\alpha_{kj}(x_k \cdot y_k)) = \alpha_k(x_k \cdot y_k),$$
so the multiplication depends neither on the particular ring $A_i$ or on the particular representation of $x$ and $y$, so $x \cdot y$ is well-defined.

Using this definition, we can easily verify other properties. Given $x,y,z \in A$, we can pass to a ring $A_i$ that has a representative for each of $x$, $y$, and $z$. Then, the representatives
satisfy the associativity $(x_iy_i)z_i = x_i(y_iz_i)$, commutativity $x_iy_i = y_ix_i$, and the distributive property $x_i(y_i + z_i) = x_iy_i = x_iz_i$, so $x$, $y$, and $z$ satisfy these properties.
Likewise, $1 \in A_i$ for all $i \in I$, and $\alpha_{ij}(1) = 1$ whenever $i \le j$. Thus, $1 = \alpha_i(1) \in A$ for any $i \in I$. Thus, if $x$ has a representative in $A_i$, $1 \cdot x = \alpha_i(1_i \cdot x_i) = \alpha_i(x_i) = x$,
so finally, $A$ is a ring.

Further, if $x,y \in A_i$, $\alpha_i(x + y) = \alpha_i(x) + \alpha_i(y)$, and by our multiplication definition, we have $\alpha_i(xy) = \alpha_i(x)\alpha_i(y)$ and $\alpha_i(1) = 1$. Thus,
$\alpha_i : A_i \to A$ is a ring homomorphism for all $i \in I$.
\end{proof}

\begin{proof}
Suppose $A_i \neq 0$ for all $i \in I$, and suppose to the contrary that $A = \varinjlim A_i = 0$. Since $A = 0$, for any $i \in I$, we have $\alpha_i(1_i) = 0$. By Exercise \ref{2.15},
there is some $j$ such that $\alpha_{ij}(1_i) = 0$. But $\alpha_{ij}(1_i) = 1_j \neq 0$ since $A_j \neq 0$. Hence, if $A = 0$, then there is some $i \in I$ such that $A_i = 0$.
\end{proof}

\subsection{Exercise}\label{2.22} % 22

\begin{proof}
First, we verify that given a directed system of ring $\{(A_i,\alpha_{ij})\}$, the nilradicals $\mathfrak{N}_i$ form a directed system $\{(\mathfrak{N}_i,\alpha_{ij})\}$. Suppose that
$n_i \in \mathfrak{N}_i$ and $j \ge i$. By definition, there is some $m$ such that $n_i^m = 0$. Thus, $\alpha_{ij}(n_i)^m = \alpha_{ij}(n_i^m) = \alpha_{ij}(0) = 0$, using that fact that
$\alpha_{ij}$ is a ring homomorphism. Thus, the directed system is well-defined.

Also, note that for each $i$, since $\mathfrak{N}_i$ is an ideal of $A_i$, we have exact sequences
$$0 \to \mathfrak{N}_i \to A_i \to A_i/\mathfrak{N}_i \to 0.$$
By Exercise \ref{2.19}, we have an exact sequence of direct limits. In particular, $\varinjlim \mathfrak{N}_i \xhookrightarrow{\iota} A$, where $A = \varinjlim A_i$. The map
$\iota = \alpha_i \circ \iota_i$, where $\iota_i : \mathfrak{N}_i \to A_i$ is just inclusion. In particular, the map $\mathfrak{N}_i \to \varinjlim \mathfrak{N}_j$ is just 
the restriction of $\alpha_i$.

If $\mathfrak{N}_A$ is the nilradical of $A$, let $n \in \mathfrak{N}_A$. By Exercise \ref{2.15}, there is some $i \in I$ and $n_i \in A_i$ such that
$\alpha_i(n_i) = n$. Since, per Exercise \ref{2.21}, $\alpha_i$ is a ring homomorphism, $\alpha_i(n_i^m) = n^m = 0$. By Exercise \ref{2.15}, there is some $j \ge i$ such that
$0 = \alpha_{ij}(n_i^m)$. But $\alpha_{ij}$ is ring homomorphism, so $n_j = \alpha_{ij}(n_i)$ satisfies $n_j^m = 0$, so $n_j \in \mathfrak{N}_j$. Further, we have 
$\alpha_j(n_j) = \alpha_j(\alpha_{ij}(n_i)) = \alpha_i(n_i) = n$. Thus, $n \in \varinjlim \mathfrak{N}_i$.

Conversely, if $n \in \varinjlim \mathfrak{N}_i$, then there is some $i \in I$ such that $n = \alpha_i(n_i)$ for some $n_i \in \mathfrak{N}_i$. Since $n_i^m = 0$ for some $m$, we have
$n^m = \alpha_i(n_i)^m = \alpha_i(n_i^m) = \alpha_i(0) = 0$, so $n \in \mathfrak{N}_A$. Thus, we have $\mathfrak{N}_A = \varinjlim \mathfrak{N}_i$.
\end{proof}

\begin{proof}
Let $A = \varinjlim A_i$, and suppose $A$ is not an integral domain. So, there are nonzero $x,y \in A$ such that $xy = 0$. By Exercise \ref{2.21}, there is some $i \in I$ such that
$x = \alpha_i(x_i)$ and $y = \alpha_i(y_i)$ for some $x_i,y_i \in A_i$. Then, 
$$0 = xy = \alpha_i(x_i)\alpha_i(y_i) = \alpha_i(x_iy_i).$$
By Exercise \ref{2.15}, there is some $j \ge i$ such that $0 = \alpha_{ij}(x_iy_i) = \alpha_{ij}(x_i)\alpha_{ij}(y_i)$ since $\alpha_{ij}$ is a ring homomorphism. Suppose $\alpha_{ij}(x_i) = 0$. Then,
$$0 = \alpha_j(0) = \alpha_j(\alpha_{ij}(x_i)) = \alpha_i(x_i) = x,$$
but we assumed $x$ is nonzero, so $\alpha_{ij}(x_i) \neq 0$. Similarly, since we assumed $y$ is nonzero, we have $\alpha_{ij}(y_i) \neq 0$. Thus, $A_j$ is not an integral domain. 
Therefore, if $A_i$ are all integral domains, then $A$ is also an integral domain.
\end{proof}

\subsection{Exercise}\label{2.23} % 23

Let $(B_{\lambda})_{\lambda \in \Lambda}$ be a family of $A$-algebras. For each finite subset of $J \subset \Lambda$, let $B_J$ denote the tensor product (over $A$) of the $B_{\lambda}$ for
$\lambda \in J$. If $J'$ is another finite subset of $\Lambda$ and $J \subset J'$, there is a canonical $A$-algebra homomorphism $\phi_{JJ'} : B_J \to B_{J'}$. Explicitly, write $B_{J'} = B_J \otimes B_K$,
where $K = J' - J$. Then, $B_J \to B_J \otimes B_K$ via $x \mapsto x \otimes \left(\bigotimes_{k \in K}1_k\right)$. Clearly, this is an $A$-algebra homomorphism. 
Note that if we have sets $I \subset J \subset K$, then the maps defined in this way satisfy $\psi_{JK}\circ\psi_{IJ} = \psi_{IK}$. Also, $\psi_{II} = \id_{B_I}$. In particular, $\{(B_I, B_{IJ})\}$ form a directed system.

Let $B$ denote the direct limit of the rings $B_J$ as $J$ runs through all finite subsets of $\Lambda$. The ring $B$ has a natural $A$-algebra structure for which the homomorphisms 
$B_J \to B$ are $A$-algebra homomorphisms (as in Exercise \ref{2.14}). 
\begin{proof}
Our canonical maps $\mu_J : B_J \to B$ are compatible $A$-module homomorphisms by the construction in Exercise \ref{2.14}. By Exercise \ref{2.21}, $\mu_J : B_J \to B$ are ring homomorphisms. Thus,
the $\mu_J$ are $A$-algebra homomorphisms.
\end{proof}
The $A$-algebra $B$ is the {\it tensor product} of the family $(B_{\lambda})_{\lambda \in \Lambda}$.

\subsection{Exercise}\label{2.24} % 24

\begin{proof}
(i) $\implies$ (ii). Suppose the $A$-module $M$ is flat, and let $N$ be an $A$-module. Then, we can take a free resolution of $N$:
$$\dots \to F_2 \to F_1 \to N \to 0.$$
Since this sequence is exact and $M$ is flat, tensoring with $M$ yields another exact sequence. Thus, the homology of this sequence is trivial,
so $\Tor_n^A(M,N) = 0$ for all $n > 0$ and all $A$-modules $N$.

(ii) $\implies$ (iii). Trivial.

(iii) $\implies$ (i). Let
$$0 \to N' \to N \to N'' \to 0$$
be any exact sequence of $A$-modules. Tensoring with $M$, we obtain the exact sequence
$$\Tor_1^A(M, N'') \to M \otimes N' \to M \otimes N \to M \otimes N'' \to 0.$$
By assumption, $\Tor_1^A(M, N'') = 0$, so the sequence is a short exact sequence. Thus, $M$ is flat.
\end{proof}

\subsection{Exercise}\label{2.25} % 25

\begin{proof}
Let $M$ be any $A$-module. Since $N''$ is flat, both $\Tor_1^A(M,N'')$ and $\Tor_2^A(M,N'')$ vanish by Exercise \ref{2.24}. Thus, tensoring
$$0 \to N' \to N \to N'' \to 0$$
with $M$, we have
$$0 \to \Tor_1^A(M,N') \to \Tor_1^A(M,N) \to 0$$
so that $\Tor_1^A(M,N') \cong \Tor_1^A(M,N)$. Thus, if one vanishes for all $M$, the other vanishes for all $M$. By Exercise \ref{2.24},
$N$ is flat if and only if $N'$ is flat.
\end{proof}

\subsection{Exercise}\label{2.26} % 26

\begin{proof}
The first direction is trivial. If $N$ is flat, then by Exercise \ref{2.24}, we have $\Tor_1(A/\mathfrak{a}, N) = 0$ for all finitely generated ideals $\mathfrak{a}$ in $A$
simply because $\Tor_1(-,N) = 0$.

For the other direction, we'll start by showing $N$ is flat if $\Tor_1(M,N) = 0$ for all finitely generated $A$-modules $M$. Suppose $M' \to M$ is an injection of finitely generated 
$A$-modules. Then,
$$\Tor_1(N, M/M') \to M' \otimes N \to M \otimes N$$
is exact. By the symmetry of $\Tor$, we have $\Tor_1(N,M/M') \cong \Tor_1(M/M',N)$. As the projection of a finitely generated $A$-module, $M/M'$ is also a finitely generated $A$-module,
so by assumption, $\Tor_1(M/M',N) = 0$. Thus, $M' \otimes N \to M \otimes N$ is injective
by exactness. Invoking Proposition 2.19, we see that $N$ is flat.

Using this, we'll show that $N$ is flat if $\Tor_1(M,N) = 0$ for all modules $M$ generated by a single element $x \in M$. By the result of the previous paragraph, its enough to show that this implies
any finitely generated module $M$ satisfies $\Tor_1(M,N) = 0$. Let $M$ be finitely generated, and let $x_1,\ldots,x_n$ be a collection of generators. Let $M_i$ be the submodule generated by $x_1,\ldots,x_i$. 
By assumption, $\Tor_1(M_1,N) = 0$. Also note that the quotients $M_i/M_{i-1}$ are generated by the projection of a single element $x_i$, so $\Tor_1(M_i/M_{i-1}) = 0$. We'll use induction
to show that $\Tor_1(M_i,N) = 0$ for all $i$ and thus that $\Tor_1(M,N) = 0$ since $M = M_n$. We'll shown the base case, so let's assume the inductive hypothesis. The exact sequence
$$0 \to M_{i-1} \to M_i \to M_i/M_{i-1} \to 0$$
gives rise to an exact sequence
$$\Tor_1(M_{i-1},N) \to \Tor_1(M_i,N) \to \Tor_1(M_i/M_{i-1},N).$$
Since both ends vanish, $\Tor_1(M_i,N) = 0$ by exactness. By induction, $\Tor_1(M,N) = 0$, so $N$ is flat whenever $\Tor_1(M,N) = 0$ for all {\it cyclic} (generated by one element) modules $M$.

We're now ready to prove the other direction. Suppose $\Tor_1(A/\mathfrak{a},N) = 0$ for all finitely generated ideals $\mathfrak{a}$. Let $M$ be any cyclic module, so $M = Ax$ for some $x \in M$.
Consider the map $\phi : A \to Ax$ defined by $a \mapsto ax$. Since $M = Ax$, this map is surjective, so we have $M \cong A/\mathfrak{a}$, where $\mathfrak{a} = \ker\phi$ is an ideal of $A$. 
Since $M \cong A/\mathfrak{a}$, then if we had $\Tor_1(A/\mathfrak{a},N) = 0$ for arbitrary ideals $\mathfrak{a}$, we would have $\Tor_1(M,N) = 0$, implying by the previous paragraph that $N$ is flat.

Thus, we just have to show that $\Tor_1(A/\mathfrak{a},N) = 0$ for all finitely generated ideals $\mathfrak{a}$ implies $\Tor_1(A/\mathfrak{b},N) = 0$ for arbitary ideals $\mathfrak{b}$. Note that
for an arbitrary ideal $\mathfrak{b}$, we have a direct system of finitely generated ideals $\mathfrak{a}$ contained in $\mathfrak{b}$, where the maps are inclusions. Since $\Tor_1(A/\mathfrak{a},N) = 0$
for all finitely generated $\mathfrak{a}$, we have exact sequences
$$0 \to \mathfrak{a} \otimes N \to A \otimes N.$$
By Exercise \ref{2.19} (direct limits preserve exact sequences) and Exercise \ref{2.20} (direct limit commutes with tensor product), we have that
$$0 \to \mathfrak{b} \otimes N \to A \otimes N$$
is exact, so $\mathfrak{b} \otimes N \to A \otimes N$ is injective. By this injectivity and the right exactness of the tensor product, we have
$$0 \to \mathfrak{b} \otimes N \to A \otimes N \to (A/\mathfrak{b}) \otimes N \to 0,$$
so $\Tor_1(A/\mathfrak{b},N) = 0$ for any ideal $\mathfrak{b}$ of $A$, completing the proof.
\end{proof}

\subsection{Exercise}\label{2.27} % 27

\begin{proof}
(i) $\implies$ (ii). Suppose $A$ is absolutely flat, and let $x \in A$. Since $A$ is absolutely flat, $(x)$ is flat, so
$$0 \to (x) \otimes (x) \to (x) \otimes A \to (x) \otimes A/(x) \to 0$$
is exact. By canonical isomorphism, this is
$$0 \to (x^2) \to (x) \xlongrightarrow{f} (x) \otimes A/(x) \to 0.$$
However, consider the commutative diagram
\begin{center}
\begin{tikzcd}
(x) \otimes A \arrow[r, "f"] \arrow[d, "\pi_1"] & (x) \otimes A/(x) \arrow[d, "\pi_2"] \\
A \arrow[r, "f'"] & A/(x)
\end{tikzcd}
\end{center}
An arbitrary generating element of $(x) \otimes A$ is $ax \otimes b$. We have 
$$f'(\pi_1(ax \otimes b)) = f'(abx) = 0$$
since $abx \in (x)$. Thus, $f' \circ \pi_1 = 0$ and thus $\pi_2 \circ f = 0$. Since $A$ is absolutely flat, $A/(x)$ is flat, so
$$0 \to (x) \otimes A/(x) \to A \otimes A/(x)$$
is injective since $(x) \to A$ is injective. But, this is just $\pi_2 : (x) \otimes A/(x) \to A/(x) \cong A \otimes A/(x)$. Thus, since $\pi_2 \circ f = 0$
and $\pi_2$ is injective, we must have $\im f = 0$. By exactness, we have that $\ker((x^2) \to (x)) = 0$, implying $(x) = (x^2)$. Thus, every principal ideal is idempotent.

(ii) $\implies$ (iii). Suppose every principal ideal is idempotent. Then, for every $x \in A$, $(x) = (x^2)$. so $x = ax^2$ for some $a \in A$. Thus, $e = ax$ is idempotent.
Note that $(e) = (x)$ since if $b \in (x)$, then $b = cx = cax^2 = cxe \in (e)$. Thus, every principal ideal is generated by a single idempotent.

Similarly, if we have a finitely generated ideal $(x_1,\ldots,x_n)$, then this is $(x_1) + \cdots + (x_n) = (e_1) + \cdots + (e_n)$, where $e_i$ is the idempotent such that
$(e_i) = (x_i)$. Thus, every finitely generated ideal is generated by idempotents. Suppose we have two idempotents $e,f \in A$. Then $(e,f) = (e + f - ef)$ since $e(e + f - ef) = e + ef - ef = e$ 
and $f(e + f - ef) = ef + f - ef = f$. But any principal ideal can be generated by an idempotent, so $(e,f) = (e+f-fg) = (g)$ for some idempotent $g$. By a routine inductive argument,
every ideal generated by finitely many idempotents is generated by one idempotent. But we saw that all finitely generated ideals have this form, so every finitely generated ideal is $(e)$ for
some idempotent $e$.

Thus, to show that every finitely generated ideal is a direct summand of $A$, we just have to show that $(e)$ is a direct summand of $A$ for every idempotent $e$. We claim that $A = (e)\oplus(1-e)$.
If $x \in A$, then $x = xe + x(1-e)$, so $(e) + (1-e) = A$. Now suppose $x \in (e) \cap (1-e)$, i.e. we have $x = ae = b(1-e)$. Then, $(a+b)e = b$. Multiplying both sides by $e$, we obtain $(a+b)e = be$. 
Thus, $b = be$, so $x = b(1-e) = b - be = 0$. Finally, $(e) \cap (1-e)$ and $A = (e) \oplus (1-e)$.

(iii) $\implies$ (i). Suppose every finitely generated ideal $\mathfrak{a}$ of $A$ is a direct summand. Let $M_{\mathfrak{a}}$ denote the complement.
Note that by injecting into the first coordinate and projecting onto the second, for any $A$-module $N$, we have the exact sequence
$$0 \to \mathfrak{a} \otimes N \to (\mathfrak{a} \otimes N) \oplus (M_{\mathfrak{a}} \otimes N) \to M_{\mathfrak{a}} \otimes N \to 0.$$
But the middle term is
$$(\mathfrak{a} \otimes N) \oplus (M_{\mathfrak{a}} \otimes N) \cong (\mathfrak{a} \oplus M_{\mathfrak{a}}) \otimes N \cong A \otimes N,$$
and $M_{\mathfrak{a}} \cong A/\mathfrak{a}$. Thus, we have an exact sequence
$$0 \to \mathfrak{a} \otimes N \to A \otimes N \to (A/\mathfrak{a}) \otimes N \to 0,$$
so $\Tor_1(A/\mathfrak{a},N) = 0$. Therefore, $N$ is flat by Exercise \ref{2.26}, and $A$ is absolutely flat.
\end{proof}

\subsection{Exercise}\label{2.28} % 28

Every Boolean ring is absolutely flat.

\begin{proof}
Every $x \in A$ satisfies $x^2 = x$. For every principal ideal $(x)$, we have $(x)^2 = (x^2) = (x)$, so by Exercise \ref{2.27}, $A$ is absolutely flat.
\end{proof}

Let $A$ be such that for any $x \in A$, there is some $n > 1$ such that $x^n = x$. $A$ is absolutely flat.

\begin{proof}
Let $(x)$ be an arbitrary principal ideal; $x$ satisfies $x^n = x$ for some $n > 1$. Thus,
$$(x) = (x^n) \subset (x^{n-1}) \subset \cdots \subset (x^2) \subset (x),$$
so all are equal. Hence, $(x) = (x^2) = (x)^2$, so by Exercise \ref{2.27}, $A$ is absolutely flat.
\end{proof}

Every homomorphic image of an absolutely flat ring is absolutely flat.

\begin{proof}
Let $A$ be absolutely flat, and let $f : A \to B$ be a homomorphism. Then, $f:A \to f(A)$ is surjective. Let $(f(x)) \in f(A)$ be an arbitrary principal ideal. By surjectivity and
the fact that $f$ is a ring homomorphism, $(f(y)) = f((y))$ for any $y \in A$. By Exercise \ref{2.27}, $(x) = (x)^2$. Thus,
$$(f(x))^2 = f((x))^2 = f((x)^2) = f((x)) = (f(x)),$$
so by Exercise \ref{2.27}, $f(A)$ is absolutely flat.
\end{proof}

If a local ring is absolutely flat, then it is a field.

\begin{proof}
Suppose $A$ is an absolutely flat local ring. Let $x \in A$ be nonzero. By Exercise \ref{2.27}, $(x) = (x)^2$ so that $x = ax^2$. Then, $e = ax$ is an idempotent. 
By Exercise \ref{1.12}, the only idempotents of a local ring are $0$ and $1$, so $ax = 1$ or $ax = 0$. If $ax = 0$, then $x = ax^2 = 0$. Since we chose $x$ to be nonzero,
we must have $ax = 1$, so $x$ is a unit. All nonzero elements are units, so $A$ is a field.
\end{proof}

If $A$ is absolutely flat, every non-unit in $A$ is a zero divisor.

\begin{proof}
Let $x \in A$ be a non-unit. By Exercise \ref{2.27}, $(x) = (x)^2$, so $x = ax^2$ for some $a \in A$. Thus, $x(1-ax) = 0$. Since $x$ is a non-unit, $1 - ax \neq 0$, so $x$
is a zero divisor.
\end{proof}

